\graphicspath{ {../../images/} }
\usetikzlibrary{external,arrows.meta}
\usepackage{caption}
\usepackage{pgfplots}
\usepackage{changepage}

\title{ITC8280 Fundamentals of Cryptography}
\subtitle{Public-key cryptography: DH, ECC \& KEMs}
\date{\today}
\author{Taaniel Kraavi}
\institute%
{%
  \textit{School of Information Technologies}\\
  \textit{Tallinn University of Technology}
}

\begin{document}
\begin{frame}
  \titlepage
\end{frame}

\begin{frame}{Modular arithmetic}
  Modular arithmetic underpins most classical public-key constructions
  \begin{itemize}[<+(1)->]
    \item \enquote{Clock arithmetic}: wrap-around behaviour
    \item Used in many hard problems:
    \begin{itemize}
      \item Combined with integer factorisation (e.g. the RSA problem)
      \item Part of algebraic structures (e.g. DLP-hard groups)
    \end{itemize}
    \item Hard problems: intractable for sufficiently large numbers
    % \begin{itemize}
    %   \item Intractable for classical computers
    %   \item Solvable (eventually) by quantum computers (Shor's algorithm)
    %   \item 48-bit factorisation claimed by Yan et al. (Dec. 2022)
    % \end{itemize}
  \end{itemize}
\end{frame}

\begin{frame}{Terminology refresher}
  Two integers are \emph{congruent} modulo $n$ if
  \[
    a \equiv b \pmod{n},
  \]
  \pause
  that is, if $n \mid a - b$ (read as: $n$ divides $a-b$).
  \pause
  We call $n$ the \alert{modulus}.

  \vspace*{1em}

  \pause
  You will often see the notation
  \[
    a \bmod n = b
  \]
  due to its familiarity among programmers:
  \begin{center}
    \texttt{b = a \% n}.
  \end{center}
\end{frame}

\begin{frame}{Sets}
  A \emph{set} is an unordered collection of distinct objects.
  \begin{itemize}[<+(1)->]
    \item $x\in S$: element $x$ is a member of $S$ (belongs to)
    \item $x\notin S$: element $x$ is not a member of $S$ (does not belong to)
    \item $\lvert S\rvert=n$: $S$ has $n$ elements (cardinality, size)
  \end{itemize}
\end{frame}

\begin{frame}{Sets}
  Important sets in this course:
  \begin{itemize}[<+(1)->]
    \item Integers without $0$: $\ZZ^*=\ZZ\setminus\{0\}=\{\dots, -2, -1, 1, 2, 3, \dots\}$
    \item Integers modulo $n$: $\ZZ_n = \{0, 1, 2, \dots, n-1\}$ {\scriptsize (sometimes denoted $\ZZ/n\ZZ$)}
    \item Nonzero integers modulo $n$: $\ZZ_n^* = \{1, 2, \dots, n-1\}$
    \item Invertible integers modulo $n$: $\ZZ_n^\times$ {\scriptsize (sometimes denoted $(\ZZ/n\ZZ)^\times$)}
    \par
    \begin{itemize}
      \item For a prime $p$, $\ZZ_p^\times = \ZZ_p^*$ 
    \end{itemize}
    \item $n$-bit bitstrings: $\{0, 1\}^n$
    \item All bitstrings: $\{0, 1\}^*$
  \end{itemize}
\end{frame}

\begin{frame}{Binary operations}
  A \emph{binary operation} $\circ$ is a mapping of elements:
  \[
    \circ: S \times S \to S
  \]
  {\scriptsize Here $\circ$ denotes an arbitrary operation (alternatively $\star, \ast, \diamond$, \dots)}

  \vspace*{1em}

  \pause
  I.e. it is a rule that for any two set elements gives another set element.
  \begin{itemize}[<+(1)->]
    \item We say that the set is \alert{closed} under the operation
    \item E.g. $\ZZ$ is closed under addition and multiplication
    \item E.g. $\NN$ is not closed under subtraction
  \end{itemize}
\end{frame}

\begin{frame}{Commutativity \& associativity}
  Let $S$ be a set and $\circ$ a binary operation on $S$.

  \pause
  $\circ$ is \emph{commutative} if for all $a, b\in S$,
  \[
      a \circ b = b \circ a\ .
  \]

  \pause
  $\circ$ is \emph{associative} if for all $a, b, c\in S$,
  \[
    a \circ (b \circ c) = (a \circ b) \circ c\ .
  \]

  \pause
  Questions:
  \begin{itemize}[<+(1)->]
    \item Addition on $\ZZ$?
    \item Subtraction on $\ZZ$?
    \item XOR $(\oplus)$ on $\{0,1\}^n$?
  \end{itemize}
\end{frame}

\begin{frame}{Groups}
  A \emph{group} is a set $\GG$ equipped with a binary operation $\circ$ on $\GG$ satisfying the following properties:
  \begin{enumerate}[<+(1)->]
    \item Associativity: for all $a, b, c\in \GG$,
    \[
      a \circ (b \circ c) = (a \circ b) \circ c
    \]
    \item Identity element: there exists an element $e\in \GG$ s.t. for all $a\in \GG$,
    \[
      a \circ e = a = e \circ a
    \]
    \item Inverse element: for each $a\in \GG$, there exists $b\in \GG$ s.t.,
    \[
      a \circ b = e = b \circ a
    \]
  \end{enumerate}
  \vspace*{-1em}
  \pause
  A group is \emph{abelian} (or commutative) if $a \circ b = b \circ a$ for all $a, b\in\GG$.
\end{frame}

\begin{frame}{Is it a group?}
  \begin{columns}[T]
    \begin{column}{0.4\textwidth}
        \begin{itemize}[<+(1)->]
        \item Is $(\NN, +)$ a group?
        \item Is $(\ZZ, +)$ a group?
      \end{itemize}
    \end{column}
    \begin{column}{0.4\textwidth}
      \begin{itemize}[<+(1)->]
        \item Is $(\ZZ^*, +)$ a group?
        \item Is $(\RR^*, \cdot)$ a group?
      \end{itemize}
    \end{column}
  \end{columns}
\end{frame}

\begin{frame}{Understanding inverses}
  Recall that division in $\RR^*$ corresponds to multiplication by the inverse:
  \[
    a \cdot b = 1 \iff b = a^{-1} = \frac{1}{a}
  \]

  \pause
  An element $a\in\ZZ_n$ has a multiplicative inverse if there exists $b\in\ZZ_n$ such that
  \[
    a\cdot b \equiv 1 \pmod{n}.
  \]
  \vspace*{-2em}
  \begin{itemize}[<+(1)->]
    \item i.e. when $a$ and $n$ are \emph{coprime} (relatively prime)
    \item We denote the set of invertible elements by $\ZZ_n^\times$
    \item E.g. $\ZZ_6^\times = \{1, 5\}$ where $1^{-1} = 1$ and $5^{-1} = 5$
    \item E.g. $\ZZ_5^\times = \{1, 2, 3, 4\}$ where  $1^{-1}=1, 2^{-1}=3, 3^{-1}=2, 4^{-1}=4$
  \end{itemize}
\end{frame}

\begin{frame}{Additive and multiplicative notation}
  \begin{center}
  \begin{tabular}{@{}l c c c c@{}}
    \toprule
    Convention & Operation & Identity & \enquote{Powers} & Inverse\\
    \cmidrule{1-5}
    Additive & $x + y$ & $0$ & $nx$ & $-x$\\
    \cmidrule{1-5}
    Multiplicative & $x\cdot y$ or $xy$ & $1$ & $x^n$ & $x^{-1}$\\
    \bottomrule
  \end{tabular}
  \end{center}

  \pause
  E.g. in the multiplicative group $\ZZ_p^*$ with $p$ prime:

  \begin{tabular}{r l}
    \pause $ab$     & $(a \cdot b) \bmod{p}$\\
    \pause $a^b$    & $a^b \bmod{p}$\\
    \pause $a^{-1}$ & $b$ s.t. $ab \equiv 1 \pmod{p}$
  \end{tabular}
\end{frame}

\begin{frame}{Cyclic groups}
  Let $\GG$ be a group.
  $g \in \GG$ \emph{generates} $\GG$ if:
  \[
    \GG = \{g^q:q\in\ZZ\} = \{\dots, g^{-1}, 1, g, g^2, g^3, \dots\}
  \]
  \pause
  We then call $g$ a generator and write $\langle g\rangle = \GG$.

  \vspace*{0.5em}

  \pause
  Example: $\ZZ_7^* = \langle 3\rangle$.
  \pause
  \begin{columns}[T]
    \begin{column}{0.3\textwidth}
      \begin{itemize}
        \item $3^0 \bmod{7} = 1$
        \item $3^1 \bmod{7} = 3$
        \item $3^2 \bmod{7} = 2$
      \end{itemize}
    \end{column}
    \begin{column}{0.3\textwidth}
      \begin{itemize}
        \item $3^3 \bmod{7} = 6$
        \item $3^4 \bmod{7} = 4$
        \item $3^5 \bmod{7} = 5$
      \end{itemize}
    \end{column}
    \begin{column}{0.2\textwidth}
    \end{column}
  \end{columns}

  \vspace*{0.5em}

  \pause
  $3^6 \bmod 7 = 1$ and the cycle repeats.

  \begin{tikzpicture}[>=Stealth,remember picture,overlay]
    \tikzset{shift={(current page.east)},xshift=-3cm,yshift=-1.05cm}

    \def\radius{3}
        
    \node (1) at (90:\radius em) {$1$};
    \node (3) at (30:\radius em) {$3$};
    \node (2) at (-30:\radius em) {$2$};
    \node (6) at (-90:\radius em) {$6$};
    \node (4) at (-150:\radius em) {$4$};
    \node (5) at (150:\radius em) {$5$};

    \draw[->] (1) -- (3);
    \draw[->] (3) -- (2);
    \draw[->] (2) -- (6);
    \draw[->] (6) -- (4);
    \draw[->] (4) -- (5);
    \draw[->] (5) -- (1);
  \end{tikzpicture}
\end{frame}

\begin{frame}{Safe primes}
  A \emph{safe} prime is a prime $p$ such that $p = 2q + 1$ with $q$ prime.
  \begin{itemize}[<+(1)->]
    \item Highly important in cryptography
    \item $\ZZ_p^*$ has $p-1$ elements
    \item $\ZZ_p^*$ has a unique subgroup of $q$ elements (subgroup of prime order)
    \item $\ZZ_p^*$ has no \enquote{useful} small subgroups
    \item Security for the discrete logarithm problem (DLP)
  \end{itemize}

  \vspace*{1em}
  \pause
  The \emph{order} of a group, denoted $\lvert\GG\rvert$, is the number of elements in the group.
\end{frame}

\begin{frame}{Discrete logarithm problem (DLP)}
  Setting:
  \begin{itemize}[<+(1)->]
    \item Let $p$ be a (large) safe prime such that $p = 2q + 1$ with $q$ prime.
    \item Fix a generator $g$ and define $\GG = \langle g \rangle = \{g^i :0 \le i < q\}$.
  \end{itemize}

  \vspace*{1em}

  \pause
  The following is then considered \emph{intractable}:
  \[
    \text{Find $x$ s.t. } g^x \equiv y \pmod{p}\ ,
  \]
  where $g, y, p$ are known.

  \pause
  No longer intractable when a crypto-relevant quantum computer (CRQC) appears!
\end{frame}

\begin{frame}{DLP security}
  The DLP exists in every finite cyclic group:
  \pause
  \begin{itemize}
    \item Not guaranteed to be hard in a \enquote{generic} group.
    \item Always research whether the DLP is intractable in a chosen group!
  \end{itemize}

  \vspace*{1em}

  \pause
  Two common group \enquote{types}:
  \begin{itemize}[<+(1)->]
    \item Group of quadratic residues modulo a safe prime
    \begin{itemize}
      \item \enquote{Standard} primes defined in \href{https://datatracker.ietf.org/doc/html/rfc3526}{RFC 3526} ($\ge3072$ bits recommended)
    \end{itemize}
    \item \enquote{Cryptographic} elliptic curve groups (slide \ref{slide:curve-ex})
  \end{itemize}

  \pause
  {\scriptsize\url{https://en.wikipedia.org/wiki/Discrete_logarithm_records}}
\end{frame}

\begin{frame}{Weak DLP group}
  Consider the group $(\ZZ_p^*, +)$.
  \begin{itemize}[<+(1)->]
    \item Not multiplicative, but great for intuition.
    \item DLP (cursed definition): $\log(y) = x \iff xg \equiv y \pmod{p}$.
    \item It follows that $x \equiv yg^{-1} \pmod{p}$.
    \item This is not hard to compute.
  \end{itemize}
\end{frame}

\begin{frame}{Public-key cryptography}
  Public-key cryptography consists of cryptographic systems that use pairs of related keys.
  \begin{itemize}[<+(1)->]
    \item Also called asymmetric cryptography
    \item Keypair: public and private/secret key
    \item Keys are linked using mathematical trapdoors (one-way functions)
    \item Public key can be known and used by anyone
    \item Private key must be kept secret
    \item $\implies$ it must be infeasible to derive the private key from the public key
  \end{itemize}
\end{frame}

\begin{frame}{Diffie-Hellman key exchange (DHKE)}
  Diffie-Hellman key exchange (DHKE) was the first known key-agreement protocol:
  \begin{itemize}[<+(1)->]
    \item Finite field DH is the \enquote{simplest}: \href{https://datatracker.ietf.org/doc/html/rfc7919}{RFC 7919} (DLP security)
    \item \enquote{Standard} groups in \href{https://datatracker.ietf.org/doc/html/rfc3526}{RFC 3526}
    \item Different variants exist, e.g. ECDH on elliptic curves
    \item Basis for more complex schemes
    \item \href{https://csrc.nist.gov/pubs/sp/800/56/a/r3/final}{NIST SP 800-56A Rev. 3}
  \end{itemize}

  \pause
  Easy to shoot yourself in the foot: a jumbled mess of choices. (my opinion)
\end{frame}

\begin{frame}{DH colours example}
  \begin{center}
    \includegraphics[height=200px]{dhex}
  \end{center}
\end{frame}

\begin{frame}{Finite field DH}
  \begin{center}
    \begin{tikzpicture}
      \node (bob)   at (0,0)     {\includegraphics[height=80px]{alice}};
      \node (alice) at (8, 0)    {\includegraphics[height=80px]{bob}};
      \node (eve)   at (4, -2.1) {\includegraphics[height=80px]{eve}};

      \node (btop) at (bob.east) [yshift=10px] {};
      \node (atop) at (alice.west) [yshift=10px] {};

      \draw[->,thick] (btop) -- (atop);

      \node (bbot) at (bob.east) [yshift=-10px] {};
      \node (abot) at (alice.west) [yshift=-10px] {};
  
      \draw[<-,thick] (bbot) -- (abot);

      \path let \p1 = (btop) in node (gxt) at (4,\y1) [yshift=7px] {$g^x$};
      \path let \p1 = (bbot) in node (gyt) at (4,\y1) [yshift=7px] {$g^y$};

      \path let \p1 = (btop) in node (circ1) at (2,\y1) {};
      \path let \p1 = (bbot) in node (circ2) at (6,\y1) {};
      
      \draw[->,thick] (circ1.center) |- (eve.west);
      \draw[->,thick] (circ2.center) |- (eve.east);

      \node[anchor=center,draw,circle,fill=white,thick,minimum size = 4px, inner sep=0pt] at (circ1) {};
      \node[anchor=center,draw,circle,fill=white,thick,minimum size = 4px, inner sep=0pt] at (circ2) {};

      \node (gx) at (eve.west) [xshift=-10px, yshift=10px] {$g^x$};
      \node (gy) at (eve.east) [xshift=10px, yshift=10px] {$g^y$};

      \node[anchor=east] (m) at (-0.5, 0.25) {$x\getsu\ZZ_q$};
      \node[anchor=east] (enc) at (-0.5, -0.25) {$g^{xy}\gets\bigl(g^y\bigr)^x$};
      \node[anchor=west] (m) at (8.8, 0.25) {$y\getsu\ZZ_q$};
      \node[anchor=west] (enc) at (8.8, -0.25) {$g^{xy}\gets\bigl(g^x\bigr)^y$};

      \node (enc) at (eve.south) [yshift=-10px] {$g^{xy}$?};
    \end{tikzpicture}
  \end{center}

  \pause
  Why can the eavesdropper not recover $x$ or $y$ and compute $g^{xy}$?
\end{frame}

\begin{frame}{DH shortcomings}
  \begin{itemize}[<+(1)->]
    \item By default, DHKE is anonymous (no authentication)
    \item Secure against sniffing/eavesdropping, but not against MITM
    \item Easy to pick a weak group (cool \href{https://weakdh.org/imperfect-forward-secrecy-ccs15.pdf}{\textit{article}})
    \item Susceptible\textsuperscript{*} to precomputation attacks: predefined groups/primes
  \end{itemize}
\end{frame}

\begin{frame}{Man-in-the-Middle attack}
  \begin{figure}
  \begin{adjustwidth}{-1em}{-1em}
    \centering
    \begin{tikzpicture}
      \node (bob)   at (0,0)     {\includegraphics[height=80px]{alice}};
      \node (alice) at (8, 0)    {\includegraphics[height=80px]{bob}};
      \node (eve)   at (4, 0) {\includegraphics[height=80px]{eve}};

      \node (btop) at (bob.east) [yshift=5px] {};
      \node (atop) at (alice.west) [yshift=5px] {};

      \node (bbot) at (bob.east) [yshift=-5px] {};
      \node (abot) at (alice.west) [yshift=-5px] {};

      \node (mltop) at (eve.west) [yshift=5px] {};
      \node (mrtop) at (eve.east) [yshift=5px] {};
      \node (mlbot) at (eve.west) [yshift=-5px] {};
      \node (mrbot) at (eve.east) [yshift=-5px] {};

      \draw[->,thick] (btop) -- node[midway, above] {$g^x$} (mltop);
      \draw[<-,thick] (bbot) -- node[midway, below] {$g^m$} (mlbot);

      \draw[->,thick] (atop) -- node[midway, above] {$g^y$} (mrtop);
      \draw[<-,thick] (abot) -- node[midway, below] {$g^m$} (mrbot);

      \node[anchor=east] (m) at (-0.5, 0.25) {$x\getsu\ZZ_q$};
      \node[anchor=east] (enc) at (-0.5, -0.25) {$g^{xm}\gets\bigl(g^m\bigr)^x$};
      \node[anchor=west] (m) at (8.8, 0.25) {$y\getsu\ZZ_q$};
      \node[anchor=west] (enc) at (8.8, -0.25) {$g^{ym}\gets\bigl(g^m\bigr)^y$};

      \node (gm) at (eve.north) [yshift=4px] {$m\getsu\ZZ_q$};
      \node (gmx) at (eve.south) [yshift=-4px] {$g^{xm} \gets \bigl(g^x\bigr)^m$};
      \node (gmy) at (gmx.south) [yshift=-4px] {$g^{ym} \gets \bigl(g^y\bigr)^m$};
    \end{tikzpicture}
  \end{adjustwidth}
  \end{figure}

  \pause
  Mallory establishes a shared secret with both Alice and Bob.
\end{frame}

\begin{frame}{Perfect forward secrecy (PFS)}
  \emph{Ephemeral} Diffie-Hellman (DHE, ECDHE):
  \begin{itemize}[<+(1)->]
    \item Server's long-term key used only for authentication / channel setup
    \item One-time \emph{session} keys generated for each session/transaction
    \item Prevents decryption of past sessions if  long-term key is compromised
    \item Prevents decryption of other sessions if session key is compromised
  \end{itemize}

  \pause
  No absolute protection against future cryptanalysis.
  \begin{itemize}
    \item PFS offers protection only as long as the algorithm remains unbroken.
  \end{itemize}

  \pause
  See also: \emph{\href{https://en.wikipedia.org/wiki/Double_Ratchet_Algorithm}{double ratchet}} for post-compromise confidentiality (e.g. Signal)
\end{frame}

\begin{frame}{Elliptic curves}
  $y^2 = x^3 + ax + b$
  \vspace*{-2em}
  \begin{center}
    \begin{tikzpicture}
    \begin{axis}[
            xmin=-5,
            xmax=5,
            ymin=-5,
            ymax=5,
            ticks=none,
            scale only axis,
            axis lines=middle,
            domain=-2.279018:3,      
            samples=201,
            smooth,   
            clip=false,
            % use same unit vectors on the axis
            axis equal image=true,
            axis line style=thick
        ]
    
    \addplot[blue,thick] {sqrt(x^3-3*x+5)} node[right] {};
    \addplot[blue,thick] {-sqrt(x^3-3*x+5)};
    \end{axis}
    \end{tikzpicture}
  \end{center}
\end{frame}

\begin{frame}{ECDL}
  Elliptic curve discrete logarithm problem:
  \begin{center}
    find $x$ given $Q = xP = \underbrace{P + P + \dots + P}_{x \text{ times}}$
  \end{center}
\end{frame}

\begin{frame}{Why elliptic curves}
  Equivalent security to \enquote{mod p groups} with better performance
  \begin{itemize}[<+(1)->]
    \item Large key-sizes for DLP-based trapdoors
    \item Best known algorithms against EC trapdoors are less efficient
    \item More efficient computations and optimisations
    \item The Internet loves them (TLS)
  \end{itemize}
\end{frame}

\begin{frame}{Why not elliptic curves}
  Complex structure
  \begin{itemize}[<+(1)->]
    \item Higher bar of expertise
    \item Supporting PK-encryption is sus (data encoding \enquote{problem})
    \item Questionable parameter generation: are there backdoors?
    \item Different side-channel concerns
    \item Patents, patents, and patents
  \end{itemize}
\end{frame}

\begin{frame}{DualEC\_DRBG}
  Did the NSA backdoor the DualEC\_DRBG ECC-based CSPRNG?
  \begin{itemize}[<+(1)->]
    \item Standardised by NIST
    \item Pushed and promoted by the NSA
    \item Thankfully retired in 2014
  \end{itemize}

  \pause
  Some links:
  \begin{itemize}
    \item {\scriptsize\url{https://blog.cryptographyengineering.com/2015/01/14/hopefully-last-post-ill-ever-write-on/}}
    \item {\scriptsize\url{https://www.schneier.com/blog/archives/2007/11/the_strange_sto.html}}
    \item {\scriptsize\url{https://youtu.be/nybVFJVXbww?si=_fn6J-TD7HnvfPtA}}
  \end{itemize}
\end{frame}

\begin{frame}{Well-known curves}
  \label{slide:curve-ex}
  Some popular curves include:
  \begin{itemize}[<+(1)->]
    \item \textbf{Curve25519} (D. J. Bernstein): good default choice
    \begin{itemize}
      \item Ed25519 EdDSA signatures
      \item X25519 DH over Curve25519
    \end{itemize}
    \item P-256, \textbf{P-384}, P-521 (NIST): FIPS-approved
    \begin{itemize}
      \item secp256r1, secp384r1, secp521r1 
    \end{itemize}
    \item Curve448: if Curve25519 is not enough
    \item \href{https://hackmd.io/@benjaminion/bls12-381}{BLS12-381}: for fancy crypto (e.g. ZK-SNARKs)
    \item Brainpool Curves (\href{https://datatracker.ietf.org/doc/html/rfc5639}{RFC 5639}): German compliance or NIST paranoia
  \end{itemize}
\end{frame}

\begin{frame}{Resources}
  Some resources:
  \begin{itemize}[<+(1)->]
    \item NIST-approved curves for key establishment: \href{https://csrc.nist.gov/pubs/sp/800/186/final}{\textit{NIST SP 800-186}}
    \item NIST-approved curves for digital signatures: \href{https://csrc.nist.gov/pubs/fips/186-5/final}{\textit{FIPS 186-5 (DSS)}}
    \item Bernstein's \href{https://safecurves.cr.yp.to}{\emph{SafeCurves}} guide (arguably biased and a bit outdated)
    \item \href{https://www.keylength.com/en/}{\textit{keylength.com}}
    \item \href{https://cryptobook.nakov.com/asymmetric-key-ciphers/elliptic-curve-cryptography-ecc}{\textit{Nakov's ECC for developers}} (a high-level explanation for devs)
    \item Jan Jancar's \href{https://neuromancer.sk/std/}{\textit{Standard curve database}}
  \end{itemize}

  \pause
  A nice \href{https://crypto.stackexchange.com/questions/93519/why-are-nist-curves-still-used}{\textit{post}} on Crypto Stack Exchange about NIST curves.
\end{frame}

\begin{frame}{Security level}
  PK cryptography typically requires larger keys to achieve the same security level as symmetric cryptosystems.

  \pause
  Some shaky explanations:
  \begin{itemize}[<+(1)->]
    \item Mathematical problems with various algorithms to attempt to solve them
    \item Relationship between public and secret key
  \end{itemize}
\end{frame}

\begin{frame}{Security level (classical adversaries!)}
  Estimates defined in \href{https://csrc.nist.gov/pubs/sp/800/57/pt1/r5/final}{NIST SP 800-57 REV. 5}

  \pause
  \begin{center}
  \begin{tabular}{@{}ccccc@{}}
    Security Strength & Symmetric & FFC & ECC & RSA\\
    \toprule
    128 & AES-128 & $L = 3072$  & $f = 256$ & $k = 3072$\\
    192 & AES-192 & $L = 7680$  & $f = 384$ & $k = 7680$\\
    256 & AES-256 & $L = 15360$ & $f = 512$ & $k = 15360$
  \end{tabular}
  \end{center}

  \pause
  Other recommendations:
  \begin{itemize}
    \item BSI: \href{https://www.bsi.bund.de/EN/Themen/Unternehmen-und-Organisationen/Standards-und-Zertifizierung/Technische-Richtlinien/TR-nach-Thema-sortiert/tr02102/tr02102_node.html}{TR-02102-1}
    \item ANSSI: \href{https://cyber.gouv.fr/le-referentiel-general-de-securite-version-20-les-documents}{RGS v2.0 annex B1}
    \item NCSC (NL): \href{https://english.ncsc.nl/publications/publications/2025/06/26/security-guidelines-for-transport-layer-security-2025-05}{IT Security Guidelines for
    Transport Layer Security (TLS)\textsuperscript{*}}
    \item {\scriptsize\url{https://www.keylength.com/en/compare/}}
  \end{itemize}
\end{frame}

\begin{frame}{Key encapsulation}
  A \emph{key encapsulation mechanism} is a triple of PPT algorithms:
  \begin{itemize}[<+(1)->]
    \item $\GEN$ is a randomised \emph{key generation algorithm}: $(\PK,\SK)\gets\GEN()$
    \item $\Encaps$ is a randomised \emph{encapsulation algorithm}: $(k, c)\gets\Encaps(\PK)$
    \item $\Decaps$ is a \emph{decapsulation algorithm}: $k\gets\Decaps(\SK, c)$.
  \end{itemize}

  \pause
  Terminology:
  \begin{itemize}[<+(1)->]
    \item $\SK$ is the secret encapsulation key, $\PK$ is its public counterpart
    \item $k$ is the shared secret (symmetric)
    \item $c$ is the encapsulated secret, also called the \emph{encapsulation}
  \end{itemize}

  \pause
  The scheme is \emph{functional} if $\Decaps(\SK, c) = k$ for $(k,c) \gets \Encaps(\PK)$.
\end{frame}

\begin{frame}{KEM scheme}
  \begin{figure}
  \begin{adjustwidth}{-1em}{-1em}
    \centering
    \begin{tikzpicture}
      \node (alice) at (0,0)      {\includegraphics[height=80px]{alice}};
      \node (bob)   at (6, 0)     {\includegraphics[height=80px]{bob}};
      \node (eve)   at (3, -2.1)  {\includegraphics[height=80px]{eve}};
  
      \draw[->,thick] (alice.east) -- (bob.west);
      
      \node[anchor=east] (gen) at (9.25, 2) {$(\PK,\SK)\gets\GEN()$};
      \draw[<-,dashed,thick] (alice.north) |- (gen.west);

      \node (el) at (eve.north) [xshift=-6px] {};
      \node (er) at (eve.north) [xshift=6px] {};

      \path let \p1 = (el) in node (circ1) at (\x1,2) {};
      \path let \p1 = (er) in node (circ2) at (\x1,0) {};

      \node[anchor=west] (enc) at (-4.25, 0) {$(k, c)\gets\Encaps(\PK)$};
      \node[anchor=east] (dec)  at (10, 0) {$k\gets\Decaps(\SK, c)$};
      
      \node (c) at (circ2.north) [yshift=5px] {$c$};

      \draw[->,dashed,thick] (circ1) -- (el);
      \draw[->,thick] (circ2.center) -- (er);

      \node[anchor=center,draw,circle,fill=white,thick,minimum size = 2px, inner sep=0pt] at (circ1) {};
      \node[anchor=center,draw,circle,fill=white,thick,minimum size = 2px, inner sep=0pt] at (circ2) {};
  
      \node (ska) at (circ1.north) [yshift=5px] {$\PK$};
    \end{tikzpicture}
  \end{adjustwidth}
  \end{figure}

  \pause
  The KEM must be IND-CCA2 secure.
\end{frame}

\begin{frame}{Hybrid encryption}
  Public-key cryptography is slow, but \enquote{solves} the key exchange problem\textsuperscript{*}.

  \pause
  Hybrid cryptosystem:
  \begin{itemize}
    \item Public-key -- \emph{key encapsulation mechanism} (KEM) or key agreement
    \item Symmetric -- \emph{data encapsulation scheme}
  \end{itemize}

  \vspace*{0.5em}

  \pause
  Typical steps:
  \begin{enumerate}[<+(1)->]
    \item Share/establish a symmetric secret using public key cryptography
    \item Derive symmetric keys from the secret with a KDF
    \item Use the symmetric keys for confidentiality + authenticity
  \end{enumerate}

  \vspace*{0.5em}

  \pause
  {\scriptsize \textsuperscript{*}Need for identity + authenticity on top of confidentiality.}
\end{frame}

\begin{frame}{Integrated Encryption Scheme (IES)}
  Informally, an \emph{integrated encryption scheme} describes how to build a hybrid encryption scheme from standard primitives.

  \vspace*{1em}

  \pause
  Required information:
  \begin{itemize}[<+(1)->]
    \item Public key mechanism
    \item Key derivation function (KDF)
    \item Symmetric encryption scheme
    \item Message authentication code system (MAC)
    \item Optional shared information
  \end{itemize}
\end{frame}

\begin{frame}{ECIES}
  The \emph{Elliptic Curve Integrated Encryption Scheme} (ECIES) is an IES instantiated with elliptic-curve Diffie-Hellman (ECDH).
  \begin{itemize}[<+(1)->]
    \item None of the relevant standards are freely accessible
    \item The construction pattern is useful regardless
    \item For elliptic curves, we use additive group notation
    \item Convention: group elements in uppercase, scalars in lowercase
  \end{itemize}
\end{frame}

\begin{frame}{ECIES approach}
  \begin{columns}[t]
    \begin{column}{0.48\textwidth}
      Encrypt $m$ knowing public key $K_B$
      \begin{enumerate}[<+(1)->]
        \item Compute $R\gets rG$ with $r\getsu\ZZ_n^*$
        \item Derive $S\gets rK_B$
        \item Derive $k_E, k_M \gets \KDF(S)$
        \item Encrypt $c \gets \ENC_k(m)$
        \item MAC $t \gets \MAC_{k_M}(c)$
        \item Output $R||c||t$
      \end{enumerate}
    \end{column}
    \begin{column}{0.48\textwidth}
      \pause
      Decrypt $R||c||t$ knowing private key $k_B$
      \begin{enumerate}[<+(1)->]
        \item Derive $S\gets k_BR$
        \item Derive $k_E, k_M \gets \KDF(S)$
        \item Abort if $t \neq \MAC_{k_M}(c)$
        \item Recover $m \gets \DEC_{k_E}(c)$
      \end{enumerate}
    \end{column}
  \end{columns}

  \vspace*{1em}

  \pause
  In practice, there are additional checks on the receiver side.
\end{frame}

\end{document}
